% RESOURCE-DIRECTSUM REPORT::SOURCE-TAILORED-NONGAUGE6::2026-10-18_1643_KST
\documentclass[11pt]{article}
\usepackage{amsmath,amssymb,amsthm,hyperref}
\newtheorem{theorem}{Theorem}
\newtheorem{lemma}[theorem]{Lemma}
\title{KPI-PRODUCT-PROBESEP6: An Adaptive Two-Probe Product Lower Bound and an Exact Lossy Promise Quotient}
\date{18 October 2026}
\begin{document}\maketitle
\noindent STATUS: PROVED/EXACT in the declared model, independent review pending; KNOWN THEOREM for KPI's lower bound; COMPUTATIONAL for finite tests. SOURCE-TAILORED-NONGAUGE6 for individual remote points, CLOCK-INDEPENDENT-GATE1, THRESHOLD-DESCENT1 and P versus NP remain OPEN.

\section{IDEAS PROPOSED BEFORE LITERATURE}
\begin{itemize}
\item[] A. Replace RESOURCE-GAUGE5's correlated graph-edge blocks by independent KPI polynomial-evaluation columns, and test joint $k$-wise independence against all alternate encodings and adaptive two-probe addresses.
\item[] B. Construct one source-aware lossy quotient by selecting a fixed basis of trace-evaluation rows. Test whether complete-block offline recovery wrongly implies two-probe online recovery.
\item[] C. Audit image-membership and uniform succinct bit access as cheap alternative computations.
\item[] D. Keep the KPI remote-point quantifier $\forall G\,\exists h\notin\mathrm{range}(G)$ separate from a hard point against every decoder.
\end{itemize}

\section{CORPUS AND LITERATURE}
Read DATA-DIRECTSUM1, PCP-GATE1, RESOURCE-GAUGE5 and its ERRATA, published Python and raw summary, UNIFORM-LOCAL-DIAG2, LOCALITY-BOUNDARY3, K-ELIMINATION-TRANSFER4 and CLOCKSAT1. The gauge result gives many individually rank-rich blocks with a single shared rank-$v-1$ table and is not a counterexample to KPI's product distribution. Primary source: Korten--Pitassi--Impagliazzo, \emph{Stronger Cell Probe Lower Bounds via Local PRGs}, ECCC TR25-030, revision of 28 May 2025, Theorem 1. It applies to arbitrary adaptive static cell-probe encoders for a public database family supporting $k$-wise-independent column distributions. Its $t=2,w=1,k\ge2\log Q$ regime gives $S=\Omega(Q/\log Q)$, subject to its full parameter hypotheses. The separate KPI remote-point theorem concerns a fixed local generator, not all alternative decoders. The 2023 Golovnev--Gur--Shinkar cell-sampling work supplies no automatic time bound for the present offline reconstruction.

\section{EXACT PRODUCT FAMILY}
Let $N=2^n$, $\mathbb F=\mathrm{GF}(2^n)$ and $Q=bN$. Each database $a=(p_1,\ldots,p_b)$ consists of independent degree-$<k$ polynomials $p_i(X)=\sum_{j<k}a_{i,j}X^j$ over $\mathbb F$. Its query bits are
\[
F((i,x),a)=\operatorname{Tr}_{\mathbb F/\mathbb F_2}(p_i(x)),\quad(i,x)\in[b]\times\mathbb F.
\]
The data structure must encode \emph{every} database and use one fixed public decoder with at most two possibly adaptive bit probes per query. Preprocessing and public query computation are free only in this declared cell-probe model.

\begin{lemma}[PRODUCT-KWISE6]
Under independent uniform coefficient vectors, any at most $k$ distinct global query bits are independent fair bits.
\end{lemma}
\begin{proof}
Within one block, evaluation at $r\le k$ distinct field points is a full-row-rank Vandermonde map, so its outputs are independent uniform field elements. The field trace is a surjective binary linear map, giving independent fair trace bits. Independent blocks preserve independence across any collection containing at most $k$ total queries.
\end{proof}

\begin{theorem}[KPI-PRODUCT-DIRECTSUM6, model-matched KPI corollary]
For $b=N^{O(1)}$, $k=\lceil C\log_2Q\rceil$, any sufficiently large fixed $C\ge2$, and sufficiently large $N$, every exact adaptive two-bit-probe representation of \emph{all} databases of this family uses
\[
\boxed{S=\Omega(Q/\log Q)=\Omega(bN/\log(bN)).}
\]
This quantifies over all source-dependent encodings, shared latent gauges, precomputed nonlinear tables, and adaptive probe addresses compatible with the same fixed public two-probe decoder.
\end{theorem}
\begin{proof}
The lemma verifies KPI Theorem 1's joint $k$-wise independence premise on all $Q$ queries. Here $t=2$ is even, $w=1$, and $2\log Q\le k\ll Q$ is inside its asymptotic regime. KPI gives $S\ge(Q/\log Q)^{2/t}(k/\log Q)^{1-2/t}t^{-1/2}2^{-O(w)}$. Substitute $t=2$. This is an instantiation of published restricted-model work, not an arbitrary-machine lower bound.
\end{proof}

\section{EXACT SOURCE-AWARE LOSSY/PROMISE QUOTIENT}
For one block, the trace truth table is $y=A_{n,k}a$ for a public binary matrix $A_{n,k}\in\mathbb F_2^{N\times kn}$. Put $d=\mathrm{rank}(A_{n,k})\le kn$. Select $d$ linearly independent rows, at public query positions $R\subset[N]$.

\begin{theorem}[PROMISED-BASIS-QUOTIENT6]
The projection $\pi_R(y)=y_R$ is bijective on the promised image $\mathrm{im}(A_{n,k})$, but each quotient value has $2^{N-d}$ preimages among unrestricted $N$-bit tables. Every promised output coordinate is reconstructible as a specified XOR of at most $d$ selected output bits. Image membership is in deterministic polynomial time by row reduction and parity checks. For $b$ product blocks, the quotient uses $bd\le bkn$ bits, with unrestricted fiber size $2^{b(N-d)}$.
\end{theorem}
\begin{proof}
The chosen rows form a basis of the row space. Thus every trace output row has a unique linear expression in the basis rows, which reconstructs the full codeword. Restriction to the image is injective and both spaces have $2^d$ elements. Outside the promise, a coordinate projection has kernel dimension $N-d$. Gaussian elimination gives the recognition algorithm.
\end{proof}

\begin{theorem}[COMPLETE-BLOCK-VERSUS-PROBE6]
In the trivial $Q$-cell one-probe truth-table representation, a single fixed set of just $bd\le bkn$ cells, independent of the database, determines \emph{all} $b$ complete blocks via source-aware offline reconstruction. Yet the KPI theorem requires $S=\Omega(bN/\log Q)$ for every online two-probe representation of the entire product family. For $k=\Theta(\log Q)$ and polynomial $b$, the storage-to-determining-set gap is at least $\Omega(N/\log^3Q)$.
\end{theorem}
\begin{proof}
Retain only the selected row-basis output locations in every block and reconstruct by the displayed parity relations. Compare $bd\le bkn$ with the KPI bound. There is no contradiction: an arbitrary reconstructed output can require $d$ online probes, not two. The set is a semantic \emph{determining set}; in DATA-DIRECTSUM1's stronger GLOBAL-BLOCK1 sense, it need not contain the original paths for all block queries.
\end{proof}

For any other valid $t$-probe representation and one fixed database, replay only the $bd$ basis queries. Their actual adaptive paths use at most $tbd$ distinct cells. Retaining their addresses and values costs $O(tbd(\log S+w))$ bits for $w$-bit cells, plus query-program and basis descriptions. Acquisition incurs $tbd$ real probes and their address-computation costs; subsequent full-block reconstruction may use up to $d$ selected-cell reads per output. The selected paths can depend on the database. This does not refute the NP-complete frozen-path optimization in RESOURCE-GAUGE5.

\section{SUCCINCTNESS, IMAGE GUARD, AND QUANTIFIERS}
For a succinct circuit $X(i,x)$ describing the $Q$ output bits, when $n,k,\log b$ are polynomial in its description length and $k=O(\log Q)$, promised product-image membership lies in PSPACE. Enumerate blocks and row indices; maintain the at most $kn$ binary row-basis vectors ($O((kn)^2)$ writable bits); obtain basis outputs by evaluating $X$; verify all parity relations on demand. The runtime may be exponential, but the workspace is polynomial. Explicit image membership is polynomial-time, so neither variant gives a new SAT-hard semantic gate.

KPI remote-point avoidance means $\forall G\,\exists h\notin\mathrm{range}(G)$. This does not imply $\exists h\,\forall G$ over arbitrary decoder programs: for any fixed $h$, a constant-output local generator hardwiring all of $h$ has zero data probes and description length $O(|h|)$. A short algorithm computing $h$ can replace hardwiring if its potentially huge CPU time is not charged. The project's UNIFORM-LOCAL-DIAG2/CLOCKSAT1 description-and-clock accounting is therefore essential and remains a separate OPEN mechanism.

\section{EXPERIMENTS, STATUS AND NEXT TARGET}
Standalone standard-library Python: experiments/meta/2026-10-18_RESOURCE_KPIPRODUCT6.py. Raw JSON: results/raw/meta/2026-10-18_RESOURCE_KPIPRODUCT6_raw.json. Exhaustive field trace, joint $k$-wise rank, promise-basis reconstruction and image-membership controls on $(n,k,b)=(3,2,2),(4,3,3),(4,4,2)$ yielded output-image ranks $d=4,5,9$, joint independent ranks $bd=8,15,18$, shared-gauge ranks $d$, and maximal selected-basis parity probe counts $3,5,7$. Every finite assertion passed. These tests do not establish the asymptotic KPI bound.

HIGHEST-VALUE NEXT TARGET: KPI-PROMISE-LOCAL-COMPILER7. Construct or falsify a semantically gated, low-succinct-complexity family whose YES side has an actual two-probe encoding and whose NO side inherits a KPI same-model product lower bound. An exact lossy quotient must preserve online probe locality (and its full description/clock ledger), not just complete-block offline recoverability. No unrestricted SAT time lower bound or P versus NP separation has been proved.

\begin{thebibliography}{2}
\bibitem{KPI} O. Korten, T. Pitassi, R. Impagliazzo, \emph{Stronger Cell Probe Lower Bounds via Local PRGs}, FOCS 2025 / ECCC TR25-030. \url{https://eccc.weizmann.ac.il/report/2025/030/}.
\bibitem{GGS} A. Golovnev, T. Gur, I. Shinkar, \emph{Derandomization of Cell Sampling}, SOSA 2023.
\end{thebibliography}
\end{document}
% END_OF_VERIFIED_REPORT::RESOURCE-DIRECTSUM::SOURCE-TAILORED-NONGAUGE6::P_VS_NP_OPEN
